\cotefile{157-2}
\scanpage{227}
\mspage{173}
% archive sheet number: 226
quel est l'axiome \unsure{mettant} en relation \unsure{le|les} \unsure{dérivateur|dérivateurs}
$\mathbb{D}$ avec $W_\infty$, ou avec quelque autre localiseur
fondamental.

Je vais faire à présent un résumé de la
théorie de $\mathrm{HOT}_W$, revue dans l'optique que je
viens de dire : \unsure{On} suppose pour $W$ \uline{les conditions}
$W(1,2,3)$ \uline{sans plus}, sauf mention explicite.

\textcircled{1} \uline{Localiseurs principaux} sur $\mathrm{Cat}_{/S}$, et sur les
\deleted{\unsure{autres}}\inserted{principales} catégories \deleted{\unsure{des diagrammes}} de $S$-catégories (cf diagramme p. \unsure{76}). Ils seront tous induits par les trois
localiseurs familiers
\diagram{\begin{tikzcd}[ampersand replacement=\&,column sep=small,row sep=small]
 W_S^d \& \& W_S^g\\
 \& W_S \arrow[ul,phantom,"\subset"{description,rotate=-45}] \arrow[ur,phantom,"\subset"{description,rotate=-135}] \&
\end{tikzcd}}
sur $\mathrm{Cat}_{/S}$. \deleted{Mais il} Il y a lieu aussi de
regarder $W_S^f$, mais jamais sur $\mathrm{Cat}_{/S}$ tout entier,
mais sur d'autres catégories \unsure{du diagramme|de diagrammes}, et
toujours dans des cas où $W_S^f$ coïncide soit avec
$W_S^d$, soit avec $W_S^g$. Mais le point de vue $W_S^f$ est
souvent plus \inserted{directement} commode, et a l'avantage de
plus d'être symétrique.

\textcircled{a} \deleted{Dans} \uline{$\mathrm{Liss}$}$_W\,S$, on a
\diagram{\begin{tikzcd}[ampersand replacement=\&,row sep=small]
 W_S^f \arrow[d,Rightarrow] \& {}=W_S^d=W_S^{\mathrm{univ}}\\
 W_S^g \arrow[d,Rightarrow] \&\\
 \underline{\hspace{1em}W_S\hspace{1em}} \& .
\end{tikzcd}}

Dualement \deleted{\unsure{dans}} \textcircled{b} \deleted{\unsure{dans}} \uline{$\mathrm{Prop}$}$_W\,S$, on a
\diagram{\begin{tikzcd}[ampersand replacement=\&,row sep=small]
 W_S^f \arrow[d,Rightarrow] \& {}=W_S^g=W_S^{\mathrm{univ}}\\
 W_S^d \arrow[d,Rightarrow] \&\\
 W_S \& {,}
\end{tikzcd}}
